\documentclass[12pt, twoside]{article}
\input{../preamble}
\usepackage{float}

\title{IB}
\author{Chris Huson}
\date{March 2026}

\fancyhead[LE]{\thepage}
%\fancyhead[RO]{\thepage \\ First \& last name: \hspace{2.25cm} \,\\ Grade: \hspace{2.25cm} \,}
\fancyhead[LO]{La Scuola d'Italia / Huson / IB Math: Sequences \\* 19 March 2026}

\fancypagestyle{firstpage}{%
  \fancyhf{}%
  \fancyhead[RO]{First \& last name: \hspace{2.25cm} \,\\ \,}%
  \fancyhead[LO]{La Scuola d'Italia / Huson / IB Math: Sequences \\* 19 March 2026}%
}

\begin{document}
\thispagestyle{firstpage}

\subsubsection*{1.2 Classwork: Introduction to Sequences}
\begin{enumerate}[itemsep=0.5cm]
\item Consider the sequence: 2, 5, 8, 11, \ldots
    \begin{enumerate}[itemsep=0.25cm]
        \item Write down the first term $u_1$. \hfill [1]
        \item Is the sequence arithmetic, geometric, or neither? \hfill [1]
        \item Find the common difference $d$. \hfill [1]
        \item Find the next two terms in the sequence. \hfill [1]
        \item Find a general expression for $u_n$, the $n^{th}$ term. \hfill [2]
    \end{enumerate}
    \begin{tikzpicture}
        \draw (0,2) rectangle (15.5,11);
        \draw [dotted] (1,10)--(14,10);
        \draw [dotted] (1,9)--(14,9);
        \draw [dotted] (1,8)--(14,8);
        \draw [dotted] (1,7)--(14,7);
        \draw [dotted] (1,6)--(14,6);
    \end{tikzpicture}

\newpage
\item Consider the sequence: 4, 12, 36, 108, \ldots
    \begin{enumerate}[itemsep=0.25cm]
        \item Write down the first term $u_1$. \hfill [1]
        \item Is the sequence arithmetic, geometric, or neither? \hfill [1]
        \item Find the common ratio $r$. \hfill [1]
        \item Find the next two terms in the sequence. \hfill [1]
        \item Find a general expression for $u_n$, the $n^{th}$ term. \hfill [2]
    \end{enumerate}
    \begin{tikzpicture}
        \draw (0,2) rectangle (15.5,11);
        \draw [dotted] (1,10)--(14,10);
        \draw [dotted] (1,9)--(14,9);
        \draw [dotted] (1,8)--(14,8);
        \draw [dotted] (1,7)--(14,7);
        \draw [dotted] (1,6)--(14,6);
    \end{tikzpicture}

\newpage
\item An arithmetic sequence has first term $u_1=5$ and common difference $d=4$.
    \begin{enumerate}[itemsep=0.25cm]
        \item Write down the first four terms of the sequence. \hfill [2]
        \item Find $u_{10}$, the tenth term. \hfill [2]
        \item Find $S_{10}$, the sum of the first 10 terms. \hfill [2]
    \end{enumerate}
    \begin{tikzpicture}
        \draw (0,2) rectangle (15.5,11);
        \draw [dotted] (1,10)--(14,10);
        \draw [dotted] (1,9)--(14,9);
        \draw [dotted] (1,8)--(14,8);
        \draw [dotted] (1,7)--(14,7);
        \draw [dotted] (1,6)--(14,6);
    \end{tikzpicture}

\newpage
\item A geometric sequence has first term $u_1=2$ and common ratio $\displaystyle r=\frac{3}{2}$.
    \begin{enumerate}[itemsep=0.25cm]
        \item Write down the first four terms of the sequence. \hfill [2]
        \item Find $u_7$, the seventh term. \hfill [2]
        \item Find $S_7$, the sum of the first 7 terms. \hfill [2]
    \end{enumerate}
    \begin{tikzpicture}
        \draw (0,2) rectangle (15.5,11);
        \draw [dotted] (1,10)--(14,10);
        \draw [dotted] (1,9)--(14,9);
        \draw [dotted] (1,8)--(14,8);
        \draw [dotted] (1,7)--(14,7);
        \draw [dotted] (1,6)--(14,6);
    \end{tikzpicture}

\end{enumerate}
\end{document}